\section{案例分析}

\subsection{\texorpdfstring{$\tau$-bench}{tau-bench}：六步完整链}

\code{TAU\_0001} 要求取消两个预订并修改第三个预订。完整源序列包含五个业务调用，加档案加载后链长为 6，在旧的 5/7 规则下会被删除。当前版本完整保留该任务。

第一步取消预订的参考调用只包含预订编号。泄露调用新增出生日期和支付方式两个可选参数，值分别来自档案中的同名字段，两条字段--工具关系均为禁止。删除这两个参数后精确恢复参考调用，后续查询、搜索和修改调用保持不变。

\subsection{BFCL：交易工作流}

\code{BFCL\_0001} 包含下单、查询订单、取消订单、查询账户和创建工单等步骤，链长为 6。下单的参考参数为订单类型、股票代码、价格和数量；泄露版本增加账户余额和支付方式。取消订单增加私密文本和交易金额。每个新增参数都对应一个档案字段。

源调用中的股票代码、价格和数量若与档案字段值相同，会对对应工具形成允许证据；账户余额等未出现在参考参数中的字段保持禁止。

\subsection{API-Bank：结构化账户对话}

\code{APIBANK\_0001} 来自一个 Level-1 对话，包含忘记密码、验证码重置、获取令牌、删除账户和重新注册五个 API 调用，链长为 6。第一次重置密码的参考调用传入状态、用户名和邮箱；泄露调用新增访问令牌和出生日期。删除账户的参考调用只传入令牌，新增邮箱和工单编号。邮箱与重置密码、邮箱与注册、访问令牌与删除账户的允许关系均有参考参数作为证据。

\subsection{AppWorld：长执行日志}

\code{APPWORLD\_0001} 要求关注满足条件的 Spotify 古典音乐艺术家，包含 15 个 API 调用。前几步读取用户信息和搜索候选艺术家；稳定抽样最终只选择第 14 步关注艺术家，在保留艺术家编号的同时额外加入邮箱和家庭地址。关注艺术家的接口不需要这两个字段。该例说明长日志不会令每个查询步骤都携带触发参数。

\subsection{Schema 冲突案例}

BFCL 一个任务调用关闭工单时传入字符串类型的工单编号，而源函数文档将其声明为整数。适配器在源边界拒绝该任务，并在报告中计为解析失败。这样构建器接收到的每个任务都满足源定义，不需要事后修补参考调用。
